\section{Non-Interactive Training}
\label{sec:noninteractive}
\subsection{Compact state and the complete update}
Refreshing every replicated diagonal separately wastes bootstrap slots. We instead reduce each gradient to one lane per feature and combine masking, scaling, and lane shifting in a plaintext product. Selected lanes from several diagonals form compact weight and momentum containers. A public compact-layout plan fixes their offsets and the required rotations; unpacking restores the periodic diagonals for the next forward pass. Bias slots can share unused compact positions. Masked packing and unpacking both have a depth cost and are included in the closed-form predictor.

The server performs encrypted forward propagation, forms the loss gradient from encrypted labels, propagates through the transposed weights, and accumulates gradients over micro-batches before one optimizer update. Increasing the effective batch beyond the physical lane count therefore amortizes optimizer refreshes without pretending the examples fit in one ciphertext. It also changes optimization behavior and cannot be treated as a cost-only substitution.

The original dense-secret bootstrap constants did not satisfy the desired precision and depth budget. The implemented preset instead uses a dense ternary main key and a temporary sparse key during modulus raising. The SSE pair is returned to the dense main key before the remaining arithmetic. Collapsed transforms and a level-optimal Chebyshev evaluator fit the selected modulus bands. This is an engineering parameter choice, not a new bootstrap construction or an inherited failure-probability theorem. Appendix~\ref{app:params} specifies the chain and the exposed key instances.

\subsection{Three refresh schedules (extension 3)}
\texttt{AfterUpdate} (AU) bootstraps every compact weight and velocity container after each update. \texttt{WhenNeeded} (WN), an additional schedule beyond the original plan, does so only if the next update would exceed the remaining depth. The decision uses the next batch's fullness: padded lanes at a linear output need a validity mask, so a partial batch can consume an extra level. An unknown next shape is conservatively treated as partial.

\texttt{AfterActivation} (AA) refreshes each hidden activation and retains the WN rule for compact weights. For the shallow eMLP-1 workload, activation refresh does not reduce the frequency of weight refresh beyond WN; it adds a bootstrap per hidden activation and micro-batch. Its motivation is enabling deeper circuits that otherwise exceed the budget, not an automatic speedup. Activation magnitudes must fit the supported bootstrap range. The observed AA subset run exceeded the conservative range in some slots and had a larger activation error; it was not clipped, and does not certify arbitrary activation distributions.

For a sequence of identical updates consuming $D$ levels with $L_a$ usable levels after unpacking, $\floor{L_a/D}$ is only an initial scheduling bound. Unpack costs between updates, partial batches, and the actual state level must also enter the scheduler. We check the structural predictor against every measured update rather than replacing this stateful schedule with a fixed average.

\subsection{Activation, optimizer, and loss (extensions 4--6)}
\paragraph{Polynomial activation.}
We retain the earlier squared activation for continuity and implement $z^2+z$ for the ReBoot comparison. A polynomial quartic option is available for depth experiments. Its additional depth and available tests do not establish a measured accuracy improvement; that planned claim remains unsupported.

\paragraph{Momentum.}
Both SGD and Nesterov-style momentum are supported. Writing $h=\eta(1+\mu)G/n$ and scaled velocity $U=\eta(1+\mu)V$ gives
\begin{align}
 U'&=\mu U+h,\nonumber\\
 W'&=W-h-\frac{\mu^2}{1+\mu}U.
 \label{eq:nag}
\end{align}
Constants are folded into exact-scale products already needed by the gradient path. A changing learning rate requires rescaling the stored velocity representation, and its lane shifts are charged. Biases use SGD; the faithful ReBoot configuration freezes their encrypted zero values while retaining the implementation's bias arithmetic. Weight decay is not implemented; the captured faithful setting uses zero decay.

\paragraph{Loss.}
A squared-error loss yields an affine output gradient; MSE normalization and ReBoot's sum reduction are matched through the learning-rate convention. The faithful configuration uses the released initialization, sample order, output padding, and schedule, not just the paper's nominal hyperparameter names. Softmax/cross-entropy approximation is optional in the plan and is not implemented or evaluated here. Block MLPs follow local classifier losses and stop gradients across blocks; they must not be described as an ordinary deeper back-propagation network.
